problem:  
**Conjecture (Cohomogeneity‑Two Rigidity for Positively Curved Manifolds).**  
Let \( M^n \) be a compact, simply connected Riemannian manifold with strictly positive sectional curvature and an effective, isometric action of a compact Lie group \( G \) of **cohomogeneity two**. Suppose the orbit space \( Q = M/G \) is homeomorphic to \( S^2 \) and carries its natural **Alexandrov metric** with curvature bounded below by \( \kappa>0 \) and identically equal to \( \kappa \) everywhere (so \( Q \) has constant Alexandrov curvature \( \kappa \)).  

Then \( M \) is conjectured to be diffeomorphic to one of the compact rank‑one symmetric spaces (CROSS):  
\[
S^n,\quad  \mathbb{C}P^k,\quad \mathbb{H}P^k,\quad \text{or } \mathrm{CaP}^2.
\]

Why is it a "good" problem:  
This conjecture forms a sharp, well‑posed rigidity question in the Grove symmetry program, testing how the geometry of the orbit space controls the topology of a positively curved manifold. The cohomogeneity‑two assumption, meaning the quotient \( Q=M/G \) is two‑dimensional, allows the use of Alexandrov geometry, where curvature can be studied via triangle comparison, geodesic angles, and singular orbit structures. Constant Alexandrov curvature on \( Q \) implies that the singular orbits meet at fixed angles determined by the local isotropy representations—a constraint expected to limit possible isometry types. Thus, the conjecture predicts that such uniform curvature downstairs rigidifies the topology upstairs to a CROSS.  

The problem connects **comparison geometry** (through Alexandrov curvature), **transformation group theory** (via isotropy and cohomogeneity analysis), and **global differential topology** (classification of positive‑curvature manifolds). Its statement is exceptionally simple—positively curved manifold, cohomogeneity two, constant‑curvature quotient—but resolving it would require deep geometric and topological insight. A proof would establish a new rigidity mechanism in positively curved geometry; a counterexample would yield the first new family of such manifolds in decades. Either outcome would have significant impact, making this a “good” research‑level problem.